\section{Limitations and Conclusion}
\label{sec:conclusion}

\subsection{Limitations}

The test result is a one-shot confirmatory evaluation on a backbone that was
frozen before the analysis; it is not a pristine blind test, because the B0 test
set had been observed in earlier project phases. The claim is therefore about a
frozen configuration, not about a new unseen benchmark.

The cut-dependent result is established on C-MAPSS with one backbone family.
Whether the crossover at 0.85--0.95 transfers to other datasets, backbones, or
degradation regimes is open. We do not claim cross-backbone generality.

The calibration baselines are deliberately simple. A stronger calibrator, or a
calibrator trained with a different objective, could change the ranking at the
late cuts. The paper's claim is about the two families we tested, not about all
possible calibrators.

The sensor-reconstruction result is limited to FD001 and FD003, three faults,
and the fixed fault-injection protocol. It is enough to remove the module from
the main method, but it is not a complete study of sensor-fault handling.

\subsection{Conclusion}

Multi-window aggregation for RUL is cut-dependent. On a frozen C-MAPSS backbone,
LTAG improves end-of-life test RMSE over a single-window baseline on all four
subsets (15/15, mean 0.391 cycles). A context-aware calibration and a linear
stack are stronger on an all-cut validation average, but their advantage is
concentrated at early cuts; at 0.85--0.95 LTAG is competitive with or better than
both. The validation mean is therefore not a reliable proxy for the end-of-life
test. A fair comparison also removes the apparent advantage of measurement-model
sensor reconstruction on bias and drift, leaving it useful only for stuck
sensors.

The practical recommendation is to report a cut-wise table alongside the
average, and to compare a learned gate against a simple calibrator before
claiming that the gate is necessary. The learned gate can still be the right
choice at the point where the maintenance decision is made.
